\documentclass[12pt]{article}
\usepackage[utf8]{inputenc}
\usepackage{amsmath}
\usepackage{amssymb}
\usepackage{graphicx}
\usepackage{booktabs}
\usepackage{listings}
\usepackage{xcolor}
\usepackage{hyperref}
\usepackage{geometry}
\geometry{margin=1in}

\lstset{
    language=Python,
    basicstyle=\ttfamily\small,
    keywordstyle=\color{blue}\bfseries,
    commentstyle=\color{green!60!black},
    stringstyle=\color{red},
    numbers=left,
    numberstyle=\tiny\color{gray},
    stepnumber=1,
    numbersep=5pt,
    backgroundcolor=\color{gray!10},
    frame=single,
    breaklines=true,
    breakatwhitespace=false,
    tabsize=4,
    showstringspaces=false
}

\title{Energy Losses Calculation Documentation}
\author{}
\date{}

\begin{document}

\maketitle

\section{Introduction}

The \texttt{energy\_losses.py} script calculates transmission line losses and converter losses for a given project configuration. It provides the foundation for line loss cost/benefit calculations and is used by both \texttt{line\_loss\_costs.py} and \texttt{emissions.py}.

\section{Method Overview}

\subsection{Purpose}

This script calculates energy losses by:
\begin{enumerate}
    \item Loading project technical and physical details
    \item Loading circuit and resistance details for the project category
    \item Calculating phase current based on capacity and system configuration
    \item Calculating full load adjustment based on line utilization
    \item Calculating line losses using resistance and current
    \item Calculating converter losses (for DC projects)
    \item Displaying results
\end{enumerate}

\subsection{Calculation Approach}

The energy loss calculation follows this structure:

\textbf{Line Losses:}
\begin{align}
I_{phase} &= \frac{P}{\sqrt{3} \times V \times N_{circuits}} \quad \text{(AC)} \\
I_{phase} &= \frac{P}{V \times N_{poles}} \quad \text{(DC)} \\
P_{loss} &= I^2 \times R \times N_{phases} \times N_{circuits} \\
\text{Losses (MWh/year)} &= P_{loss} \times \text{Full Load Adj} \times 8760
\end{align}

\textbf{Converter Losses (DC only):}
\begin{align}
P_{converter} &= P \times \text{Utilization} \times \text{Loss Rate} \\
\text{Converter Losses (MWh/year)} &= P_{converter} \times 8760
\end{align}

\subsection{Inputs and Outputs}

\textbf{Inputs:}
\begin{itemize}
    \item Project technical details (construction type, AC/DC, capacity, conductor type, converter type, utilization)
    \item Physical details (line length)
    \item Circuit and resistance details (voltage, conductors per phase, number of phases, circuits/poles, resistance values)
\end{itemize}

\textbf{Outputs:}
\begin{itemize}
    \item Line losses (MW/mile, MWh/year, percentage)
    \item Converter losses (MW, MWh/year, percentage)
    \item Total losses (line + converter)
    \item Lifetime losses
\end{itemize}

\subsection{Integration with Other Scripts}

This script integrates with:
\begin{itemize}
    \item \texttt{line\_loss\_costs.py} - Uses energy losses to calculate costs/benefits (see \texttt{line\_loss\_costs\_documentation.tex})
    \item \texttt{emissions.py} - Uses energy losses to calculate emissions costs (see \texttt{emissions\_documentation.tex})
    \item \texttt{yaml\_loaders.py} - Loads project configuration (see \texttt{utilities\_documentation.tex})
    \item \texttt{calculation\_utils.py} - Provides calculation functions (see \texttt{utilities\_documentation.tex})
\end{itemize}

\section{Key Functions}

\subsection{\texttt{print\_results(...)}}

Prints organized results in sections (system configuration, line losses, converter losses, summary).

\textbf{Parameters:}
\begin{itemize}
    \item Multiple parameters including voltage, phase current, resistance, losses, etc.
\end{itemize}

\subsection{\texttt{main()}}

Main execution function that orchestrates the energy loss calculation.

\textbf{Workflow:}
\begin{enumerate}
    \item Loads project technical details and constructs category identifier
    \item Determines number of converters (0 for AC projects)
    \item Loads physical details (line length)
    \item Loads circuit and resistance details
    \item Calculates phase current using \texttt{calculate\_phase\_current()}
    \item Calculates full load adjustment using \texttt{full\_load\_adjusted()}
    \item Calculates line losses using \texttt{calculate\_line\_losses()}
    \item Calculates converter losses using \texttt{calculate\_converter\_losses()} (DC only)
    \item Displays results
\end{enumerate}

\textbf{Key Logic:}
\begin{lstlisting}
# Construct category identifier
category = f"{construction_type}/{ac_dc}/{capacity_mw}MW/{conductor_type}/{converter_type}"

# Load circuit and resistance details
voltage_kv, conductors_per_phase, number_of_phases, 
number_of_circuits_poles, AC_75_resistance, DC_20_resistance = (
    load_circuit_and_resistance_details(category)
)

# Calculate phase current
phase_current = calculate_phase_current(
    capacity_mw, voltage_kv, number_of_phases, 
    number_of_circuits_poles, ac_dc
)

# Calculate full load adjustment
full_load_adj = full_load_adjusted(line_utilization_percent)

# Calculate line losses
losses_mwh_per_year, lifetime_losses_mwh, ... = calculate_line_losses(
    phase_current, full_load_adj, AC_75_resistance, DC_20_resistance,
    ac_dc, number_of_circuits_poles, conductors_per_phase,
    number_of_phases, line_length, project_lifetime,
    capacity_mw_numeric, line_utilization_percent
)

# Calculate converter losses (DC only)
if ac_dc == "DC":
    total_converter_losses_mw, total_converter_losses_mwh, ... = (
        calculate_converter_losses(
            number_of_converters, converter_type,
            line_utilization_percent, capacity_mw_numeric, ac_dc
        )
    )
\end{lstlisting}

\section{Example}

The following example demonstrates a typical usage:

\begin{lstlisting}
# Example: 500 MW Overhead AC line
# Category: "Overhead/AC/500MW/Advanced Aluminum/NA"
# 
# Inputs (from YAML):
#   - Capacity: 500 MW
#   - Voltage: 345 kV
#   - Conductors per phase: 2
#   - Number of phases: 3
#   - Number of circuits: 1
#   - AC resistance (75°C): 0.05 ohms/mile
#   - Line utilization: 70%
#   - Line length: 100 miles
#   - Project lifetime: 50 years
#
# Calculation:
#   Phase current = 500 MW / (√3 × 345 kV × 1) = 837 A
#   Full load adj = 0.7² = 0.49
#   Line loss = (837² × 0.05 × 3 × 1) × 0.49 = 51.4 MW
#   Line loss % = 51.4 / 500 = 10.3%
#   Losses MWh/year = 51.4 × 8760 = 450,000 MWh/year
#   Lifetime losses = 450,000 × 50 = 22.5M MWh
#
# Output:
#   - Line loss: 0.514 MW/mile
#   - Line loss %: 10.3%
#   - Losses MWh/year: 450,000
#   - Lifetime losses: 22.5M MWh
\end{lstlisting}

\textbf{Integration Example:}

\begin{lstlisting}
# Used by line_loss_costs.py:
from energy_losses import (
    load_project_technical_details,
    load_physical_details,
    load_circuit_and_resistance_details,
    calculate_phase_current,
    full_load_adjusted,
    calculate_line_losses,
)

# Calculate losses for a configuration
losses_mwh_per_year, lifetime_losses_mwh = calculate_configuration_losses(
    construction_type, ac_dc, capacity_mw, conductor_type,
    converter_type, line_utilization_percent, project_lifetime
)
\end{lstlisting}

\section{Key Implementation Details}

\subsection{AC vs. DC Calculations}

The script handles both AC and DC systems:
\begin{itemize}
    \item \textbf{AC}: Uses $\sqrt{3}$ factor, resistance at 75°C
    \item \textbf{DC}: No $\sqrt{3}$ factor, resistance at 20°C
\end{itemize}

\subsection{Full Load Adjustment}

The full load adjustment accounts for the fact that losses are proportional to current squared, and average current is less than peak current due to utilization. The adjustment is typically $u^2$ where $u$ is the utilization factor.

\subsection{Converter Losses}

Converter losses only apply to DC projects and are calculated based on:
\begin{itemize}
    \item Number of converters
    \item Converter type (loss rate)
    \item Line utilization
    \item Capacity
\end{itemize}

\subsection{Resistance Values}

The script uses different resistance values:
\begin{itemize}
    \item \textbf{AC}: Resistance at 75°C (operating temperature)
    \item \textbf{DC}: Resistance at 20°C (reference temperature)
\end{itemize}

\subsection{Multiple Circuits/Poles}

The script accounts for multiple circuits (AC) or poles (DC) by multiplying losses by the number of circuits/poles.

\subsection{Standalone vs. Library Use}

This script can be run standalone (displays results) or imported as a library (used by other scripts). The calculation functions are in \texttt{calculation\_utils.py} for reuse.

\end{document}

